\section{我们到底在评估什么}

最后一个核心问题是：你到底在评估 model、system，还是 method？

\subsection{方法还是系统}

在早期机器学习里，人们更多评估的是 \textbf{方法}：给定标准数据集和标准训练流程，谁的算法更好。今天很多场景里，我们评估的是 \textbf{模型/系统}：只要最后用户体验更好，内部用了什么 prompt、工具或 scaffold 并不重要。

\subsection{为什么这很重要}

\begin{itemize}
    \item 如果你在做研究，可能更关心方法是否普适、是否能推广。
    \item 如果你在做产品，可能只关心系统最终表现。
    \item 如果规则不明确，不同团队就会在不同层面“赢得”同一个 benchmark。
\end{itemize}

\subsection{几个例外}

课程也提到了一些仍然评估方法的例子：

\begin{itemize}
    \item \textbf{nanoGPT speedrun}：固定数据，比较谁更快达到某个验证损失。
    \item \textbf{DataComp-LM}：给定原始数据，按标准训练管线比较最终效果。
\end{itemize}

\begin{importantbox}{规则的游戏必须说清楚}
一个好的评估一定要先声明：输入是什么、允许什么工具、是否允许额外训练、是测最终系统还是测单模型。否则分数本身没有可比性。
\end{importantbox}

